\documentclass[11pt]{article}
\usepackage[margin=1in]{geometry}
\usepackage{amsmath,amssymb,amsthm,booktabs,longtable,hyperref,array}
\usepackage[T1]{fontenc}
\usepackage{stix2}
\usepackage{microtype}
\usepackage{enumitem}
\usepackage{mathtools}
\hypersetup{colorlinks=true,linkcolor=blue,citecolor=blue,urlcolor=blue,unicode=true}

\newtheorem{definition}{Definition}[section]
\newtheorem{lemma}[definition]{Lemma}
\newtheorem{theorem}[definition]{Theorem}
\newtheorem{proposition}[definition]{Proposition}
\newtheorem{corollary}[definition]{Corollary}
\newtheorem{remark}[definition]{Remark}
\newtheorem{conjecture}[definition]{Conjecture}

\title{A Ternary $(4k+1(2))/3$ Collatz-Type Map and Its $4^r$-Scaled Family}
\author{Farhad Banazadeh\\\texttt{fn.bana@hotmail.com}\\ORCID: 0009-0004-7023-0298}
\date{6 September 2026}

\begin{document}
\maketitle

\begin{abstract}
This paper develops a scaled family for the ternary $(4k+1(2))/3$ Collatz-type map studied in the author's preceding computational work. The base system starts from positive integers not divisible by $3$, applies the pre-reduction rule $G(3p+1)=4p+2$ and $G(3p+2)=4p+3$, and then removes all factors of $3$. The compressed base map has the two observed cycles $C_1=(1,2)$ and $C_7=(7,10,14,19,26,35,47)$. For every integer $r\ge0$, we introduce the generalized pre-reduction map
\[
M_r(x)=4x+4^r\bigl(3-(x\bmod 3)\bigr),
\]
with branches $4x+2\,4^r$ for $x\equiv1\pmod3$ and $4x+4^r$ for $x\equiv2\pmod3$, and then apply the same complete $3$-adic compression. On the scaled domain $S_r=4^rD$, where $D=\{n\in\mathbb N:3\nmid n\}$, the resulting map $T_r$ is exactly conjugate to the base map through $S_r(y)=4^ry$:
\[
T_r(4^ry)=4^rT_0(y).
\]
The identity is proved from the branch formulas and the invariance of the $3$-adic valuation under multiplication by $4^r$. An induction gives $T_r^j(4^ry)=4^rT_0^j(y)$ for every $j\ge0$. Consequently, cycles, cycle lengths, stopping times, basin membership, and all orbit relations are transported exactly by scaling, while orbit peak values are multiplied by $4^r$. In particular, the two observed base cycles become $(4^r,2\,4^r)$ and $4^r\{7,10,14,19,26,35,47\}$. The large finite computation from the base paper, which exhaustively tested all admissible starts through $10^9$, therefore transfers algebraically to every scaled level on the corresponding scaled domain. The global convergence statement remains a conjecture: the scaling theorem is exact, but it does not by itself prove convergence of the base system for all positive integers.
\end{abstract}

\section{Introduction and problem formulation}
The classical Collatz problem shows how a simple integer rule can produce difficult questions about long-term dynamics. Many Collatz-type systems change the multiplier, the divisor, the residue classes, or the amount of compression. The present work concerns a ternary system built from the two residue classes modulo $3$ and complete removal of factors of $3$ after each pre-reduced step.

The base map studied in the author's preceding paper is defined on positive integers not divisible by $3$. If $n=3p+1$ then the pre-reduced value is $4p+2$; if $n=3p+2$ it is $4p+3$. All factors of $3$ are then removed. The resulting compressed system has two observed cycles, $C_1$ and $C_7$, and was exhaustively checked for every admissible starting value up to $10^9$.

The purpose here is different. We ask whether the constants in the positive branch can be enlarged by powers of $4$ without changing the underlying dynamics. The answer is exact: the powers of $4$ produce a linearly conjugate family. Thus the new systems are not new dynamics in the sense of creating new cycle lengths or new stopping-time behavior; they are scaled copies of the base system.

The paper follows the structure of the supplied theoretical notes: formal definitions, well-definedness, $3$-adic valuation, the scaling conjugacy, induction for iterates, cycle preservation, conditional convergence, explicit cycles, stopping-time and peak consequences, and computational interpretation. The numerical data used below for the base system are kept separate from the new scaled-family statements.

\section{The base ternary map}
Let
\[
D=\{n\in\mathbb N:3\nmid n\}.
\]
Every $n\in D$ has exactly one representation $n=3p+1$ or $n=3p+2$ with $p\ge0$.

\begin{definition}[Base pre-reduction]
Define $G:D\to\mathbb N$ by
\[
G(3p+1)=4p+2,\qquad G(3p+2)=4p+3.
\]
Equivalently,
\[
G(n)=\left\lfloor\frac{4n+2}{3}\right\rfloor.
\]
\end{definition}

For a positive integer $m$, define
\[
v_3(m)=\max\{j\ge0:3^j\mid m\}.
\]

\begin{definition}[Base compressed map]
The base map is
\[
T_0(n)=\frac{G(n)}{3^{v_3(G(n))}},\qquad n\in D.
\]
Thus the output is again not divisible by $3$.
\end{definition}

The computational convention in the base study identifies the eliminated state $3$ with $1$. Hence
\[
1\longrightarrow2\longrightarrow1.
\]
The central seven-cycle is
\[
7\to10\to14\to19\to26\to35\to47\to7.
\]
Indeed, $G(47)=63=3^2\cdot7$, so the last compressed step is $47\to7$.

\section{The $4^r$-scaled family}
Fix an integer $r\ge0$. Since $4\equiv1\pmod3$, multiplication by $4^r$ preserves both nonzero residue classes modulo $3$.

\begin{definition}[Scaled domain]
Define
\[
S_r=4^rD=\{4^ry:y\in D\}.
\]
Every $x\in S_r$ has a unique representation $x=4^ry$ with $y\in D$.
\end{definition}

\begin{definition}[Scaled pre-reduction]
For $x\in S_r$, define
\[
M_r(x)=4x+4^r\bigl(3-(x\bmod3)\bigr).
\]
Equivalently,
\[
M_r(x)=
\begin{cases}
4x+2\,4^r,&x\equiv1\pmod3,\\
4x+4^r,&x\equiv2\pmod3.
\end{cases}
\]
\end{definition}

The notation keeps $k$ for the parameter in the base expressions $4k+1$ and $4k+2$, while $n$ or $x$ denotes a generic state. This avoids replacing the branch parameter by the state variable without explanation.

\begin{lemma}[Well-definedness]
For every $r\ge0$ and every $x\in S_r$, $M_r(x)$ is divisible by $3$.
\end{lemma}
\begin{proof}
Modulo $3$, we have $4^r\equiv1$ and $4\equiv1$. Therefore
\[
M_r(x)\equiv x+3-(x\bmod3)\equiv0\pmod3.
\]
More explicitly, if $x\equiv1$ then $M_r(x)\equiv4x+2\equiv1+2\equiv0$, while if $x\equiv2$ then $M_r(x)\equiv4x+1\equiv2+1\equiv0$.
\end{proof}

\begin{definition}[Scaled compressed map]
For $x\in S_r$, define
\[
T_r(x)=\frac{M_r(x)}{3^{v_3(M_r(x))}}.
\]
The image again belongs to $S_r$.
\end{definition}

\section{The key $3$-adic scaling lemma}
The exact relation between the scaled and base pre-reduction maps is the central algebraic fact.

\begin{lemma}[Scaling identity]
For every $r\ge0$ and every $y\in D$,
\[
M_r(4^ry)=4^rM_0(y).
\]
Moreover,
\[
v_3\bigl(M_r(4^ry)\bigr)=v_3\bigl(M_0(y)\bigr).
\]
\end{lemma}
\begin{proof}
Because $4^r\equiv1\pmod3$, we have $4^ry\equiv y\pmod3$, so both arguments use the same branch. Using the unified formula,
\[
\begin{aligned}
M_r(4^ry)
&=4(4^ry)+4^r\bigl(3-(y\bmod3)\bigr)\\
&=4^r\Bigl(4y+3-(y\bmod3)\Bigr)\\
&=4^rM_0(y).
\end{aligned}
\]
Since $\gcd(4^r,3)=1$, multiplication by $4^r$ does not change the exponent of $3$ in an integer. Hence
\[
v_3(4^rM_0(y))=v_3(M_0(y)).
\]
\end{proof}

\section{Exact scaling conjugacy}
\begin{theorem}[One-step conjugacy]
For every $r\ge0$ and every $y\in D$,
\[
T_r(4^ry)=4^rT_0(y).
\]
Equivalently, with $S_r(y)=4^ry$,
\[
T_r\circ S_r=S_r\circ T_0,
\]
and therefore
\[
T_r=S_r\circ T_0\circ S_r^{-1}
\]
on the scaled domain $S_r$.
\end{theorem}
\begin{proof}
By the preceding lemma,
\[
M_r(4^ry)=4^rM_0(y)
\]
and the two quantities have the same $3$-adic valuation. Therefore
\[
T_r(4^ry)
=\frac{4^rM_0(y)}{3^{v_3(M_0(y))}}
=4^rT_0(y).
\]
\end{proof}

\begin{theorem}[Conjugacy of all iterates]
For every $j\ge0$, every $r\ge0$, and every $y\in D$,
\[
T_r^j(4^ry)=4^rT_0^j(y).
\]
\end{theorem}
\begin{proof}
Induction on $j$. For $j=0$ both sides equal $4^ry$. Assume the identity for $j$. Then
\[
\begin{aligned}
T_r^{j+1}(4^ry)
&=T_r\bigl(T_r^j(4^ry)\bigr)\\
&=T_r\bigl(4^rT_0^j(y)\bigr)\\
&=4^rT_0\bigl(T_0^j(y)\bigr)\\
&=4^rT_0^{j+1}(y),
\end{aligned}
\]
where the third equality is the one-step conjugacy theorem.\end{proof}

\section{Cycle preservation}
\begin{theorem}[Exact preservation of cycles]
The cycles of $T_r$ on $S_r$ are exactly the $4^r$-scaled cycles of $T_0$, with the same lengths.
\end{theorem}
\begin{proof}
If $C=(c_0,\ldots,c_{L-1})$ is a cycle of $T_0$, then
\[
T_0^L(c_0)=c_0.
\]
By the iterate identity,
\[
T_r^L(4^rc_0)=4^rT_0^L(c_0)=4^rc_0.
\]
Thus $4^rC$ is a cycle of the same length. Conversely, if $x=4^ry\in S_r$ is periodic under $T_r$, then for some $L\ge1$,
\[
4^ry=T_r^L(4^ry)=4^rT_0^L(y),
\]
so $T_0^L(y)=y$. Hence $y$ belongs to a base cycle. No additional cycles exist in $S_r$.
\end{proof}

The two observed cycles therefore become
\[
C_1^{(r)}=(4^r,\,2\,4^r),
\]
and
\[
C_7^{(r)}=4^r(7,10,14,19,26,35,47).
\]
The second cycle closes because
\[
M_r(47\,4^r)=4^rM_0(47)=4^r\cdot63,
\]
so the compressed output is
\[
\frac{4^r\cdot63}{3^2}=7\,4^r.
\]

\section{Stopping times, peaks, and basin membership}
Let the base orbit of $y$ be $y_j=T_0^j(y)$ and define the first-descent stopping time, when it exists, by
\[
\sigma_0(y)=\min\{j\ge1:T_0^j(y)<y\}.
\]
For the scaled orbit $x_j=T_r^j(4^ry)$, the iterate theorem gives
\[
x_j=4^ry_j.
\]
Therefore
\[
x_j<4^ry\iff y_j<y,
\]
and hence
\[
\sigma_r(4^ry)=\sigma_0(y).
\]
Thus stopping times are invariant under scaling.

If
\[
P_0(y)=\max_{0\le j\le\sigma_0(y)}y_j,
\]
then the corresponding scaled peak is
\[
P_r(4^ry)=4^rP_0(y).
\]
The same scaling holds for any finite orbit segment and, in particular, for a global peak measured over a matched trajectory segment.

Basin membership is also preserved exactly:
\[
4^ry\in\mathcal B_r(4^rC)\iff y\in\mathcal B_0(C).
\]
Consequently the number of starts in a scaled basin is unchanged only when the compared base and scaled domains contain corresponding sets of starts. For a fixed numerical cutoff $N$, the relevant base cutoff is $\lfloor N/4^r\rfloor$, so basin counts at a fixed $N$ should not be claimed to be identical.

\section{The $3$-adic and average-logarithmic viewpoint}
The base pre-reduced value has the large-scale size
\[
G(n)\approx\frac43n.
\]
If $v=v_3(G(n))$, then complete compression gives the approximate ratio
\[
\frac{T_0(n)}n\approx\frac{4}{3^{v+1}}.
\]
Under the Haar-measure or uniform-residue heuristic in $\mathbb Z_3$,
\[
\Pr(v=k)=\frac{2}{3^{k+1}},\qquad k\ge0,
\]
and therefore
\[
E[v]=\frac12.
\]
The corresponding average logarithmic drift is
\[
\delta\approx\ln\frac43-\frac12\ln3
=\ln\left(\frac{4}{3\sqrt3}\right)
\approx-0.2616.
\]
The associated geometric mean factor is
\[
\rho=\exp(\delta)=\frac{4}{3\sqrt3}\approx0.7698,
\]
which corresponds to an average logarithmic contraction of about $23.02\%$ in this heuristic model.

The same heuristic scale is inherited by the scaled family because
\[
\frac{T_r(4^ry)}{4^ry}=\frac{T_0(y)}y.
\]
Thus the scaling parameter $r$ changes the absolute size of states but not the normalized one-step ratio along corresponding trajectories.

This argument is deliberately a heuristic. The residue classes along a deterministic orbit are not independent random samples, so negative average drift under a Haar model is not a proof that every orbit converges.

\section{Explicit cycle calculations}
The two observed base cycles are
\[
C_1=(1,2)
\]
and
\[
C_7=(7,10,14,19,26,35,47).
\]
For $C_1$,
\[
M_0(1)=6=2\cdot3,\qquad T_0(1)=2,
\]
and
\[
M_0(2)=9=3^2,\qquad T_0(2)=1.
\]
Hence the scaled two-cycle is
\[
4^r\to2\,4^r\to4^r.
\]

For the seven-cycle,
\[
7\to10\to14\to19\to26\to35\to47,
\]
with no $3$-adic compression on the first six transitions. At the last state,
\[
M_0(47)=63=3^2\cdot7,
\]
so $47\to7$. Multiplying every state by $4^r$ gives
\[
7\,4^r\to10\,4^r\to14\,4^r\to19\,4^r\to26\,4^r\to35\,4^r\to47\,4^r\to7\,4^r.
\]
For example, at $r=1$ this is
\[
28\to40\to56\to76\to104\to140\to188\to28,
\]
and the two-cycle is
\[
4\to8\to4.
\]
At $r=2$ the seven-cycle is
\[
112\to160\to224\to304\to416\to560\to752\to112,
\]
and the two-cycle is
\[
16\to32\to16.
\]
These are exact algebraic consequences of the conjugacy theorem, not imported computational data.

\section{Computational evidence inherited from the base system}
The preceding base paper exhaustively tested every admissible start $1\le n\le10^9$ with $3\nmid n$. The exact reported finite result was:

\begin{table}[h]
\centering
\caption{Base-system exhaustive computation through $10^9$.}
\begin{tabular}{lr}
\toprule
Quantity & Value\\
\midrule
Upper bound $N$ & 1,000,000,000\\
Admissible starts & 666,666,667\\
$C_1$ basin & 21,785,111\\
$C_1$ percentage & 3.267766648366117\%\\
$C_7$ basin & 644,881,556\\
$C_7$ percentage & 96.732233351633880\%\\
Unknown & 0\\
Total steps & 42,759,887,447\\
Mean steps & 64.139831138430083\\
Maximum steps & 385\\
Start attaining maximum & 696,171,200\\
Global peak & 134,701,251,885,711,310\\
Start attaining global peak & 920,435,228\\
Global-peak step & 115\\
Total removed $3$-factors & 21,683,837,513\\
Maximum removed in one step & 19\\
Runtime & 527.834688 s\\
\bottomrule
\end{tabular}
\end{table}

These values belong to the base computation and are not a new billion-scale computation of the generalized family. The scaling theorem supplies the exact correspondence: every tested base trajectory starting at $y$ has a scaled trajectory starting at $4^ry$ with the same number of steps and the same cycle destination after multiplication by $4^r$. Peak values are multiplied by $4^r$. Therefore a computational run for the base system can be reused algebraically on every scaled copy without repeating the dynamics.

For a fixed scaled level $r$, the set corresponding to all base starts $y\le N$ is $4^rD\cap[1,4^rN]$. Thus the complete base computation through $10^9$ immediately gives a complete corresponding scaled-domain statement through $4^r10^9$.

\section{Conjecture for the scaled family}
\begin{conjecture}[Two-cycle convergence for the scaled family]
For every $r\ge0$ and every $y\in D$, the orbit of $4^ry$ under $T_r$ eventually enters one of
\[
C_1^{(r)}=(4^r,2\,4^r)
\]
and
\[
C_7^{(r)}=4^r(7,10,14,19,26,35,47).
\]
Equivalently, the base orbit of every positive integer under $T_0$ eventually enters $C_1$ or $C_7$.
\end{conjecture}

\begin{theorem}[Conditional convergence transfer]
If every base orbit in $D$ converges to one of $C_1$ or $C_7$, then every orbit of $T_r$ on $S_r$ converges to the corresponding scaled cycle for every $r\ge0$.
\end{theorem}
\begin{proof}
Let $x=4^ry$. By the iterate identity,
\[
T_r^j(x)=4^rT_0^j(y).
\]
If the base orbit eventually enters $C_i$, then after the same number of steps the scaled orbit enters $4^rC_i$.\end{proof}

The converse is also immediate: if every scaled system converges, taking $r=0$ gives convergence of the base system. Hence the scaled-family convergence problem contains no independent global difficulty beyond the base conjecture.

\section{What the scaling theorem proves and what it does not prove}
The exact results of this paper are algebraic. They prove:
\begin{enumerate}[label=(\roman*)]
\item the scaled pre-reduction is well-defined;
\item multiplication by $4^r$ preserves the relevant residue class modulo $3$;
\item the $3$-adic valuation is unchanged by the scale factor;
\item the one-step maps are exactly conjugate;
\item all iterates correspond exactly;
\item cycles and their lengths correspond exactly;
\item stopping times and normalized orbit ratios are preserved;
\item peak values scale by $4^r$;
\item basin membership transfers exactly on corresponding scaled domains.
\end{enumerate}

What the theorem does not prove is global convergence of the base system. The finite computation through $10^9$ is complete for that finite domain, but it cannot rule out a different cycle or a divergent orbit entirely above the tested range. The negative logarithmic drift is a heuristic, not a deterministic proof.

\section{Relation to the earlier theoretical note}
The supplied theoretical note studies the positive family with constants $2\,4^r$ and $4^r$ in the two branches and defines
\[
M_r(x)=
\begin{cases}
4x+2\,4^r,&x\equiv1\pmod3,\\
4x+4^r,&x\equiv2\pmod3,
\end{cases}
\]
followed by division by the full power of $3$. It establishes the scaling identity $M_r(4^ry)=4^rM_0(y)$, equality of the $3$-adic valuations, one-step conjugacy, induction for iterates, and exact preservation of cycles. The present paper incorporates those proofs into the notation of the actual compressed ternary map and connects them directly to the base computational record.

The model material is used here as a structural and theoretical template. Its numerical experiments are not used as numerical evidence for the present paper. In particular, no finite-range counts, runtimes, stopping-time records, or other numerical values from the model article are merged with the base $10^9$ computation. The only numerical examples in the scaled-family sections are exact consequences of the formulas or values explicitly inherited from the base paper.

\section{Reproducibility and data separation}
The supplementary archive accompanying this article should contain:
\begin{itemize}
\item this LaTeX source;
\item the compiled PDF;
\item a small exact verifier for the conjugacy identity;
\item a machine-readable statement of the scaling rules;
\item a README explaining which results are inherited from the base paper and which are proved here;
\item the base-paper numerical summary as a clearly labelled reference, not as a new computation.
\end{itemize}

The base computational record is available in the author's previous Zenodo record, DOI 10.5281/ZENODO.22195651, together with the associated source and data. The present archive is intended to make the scaled-family derivation independently reproducible.

\section{Conclusion}
A natural $4^r$ generalization of the ternary $(4k+1(2))/3$ compressed map has been constructed and analyzed. The generalized pre-reduction rule is
\[
M_r(x)=4x+4^r\bigl(3-(x\bmod3)\bigr),
\]
with branches $4x+2\,4^r$ and $4x+4^r$. On the scaled domain $S_r=4^rD$, this family is exactly conjugate to the base map:
\[
T_r(4^ry)=4^rT_0(y).
\]
The proof depends on two simple facts: $4^r\equiv1\pmod3$ and $\gcd(4^r,3)=1$. The first preserves the branch, and the second preserves the $3$-adic valuation used in compression.

The consequences are exact. Every orbit is a scaled copy of a base orbit, every cycle is scaled without changing its length, stopping times are unchanged, and peaks are multiplied by $4^r$. The observed cycles become $(4^r,2\,4^r)$ and $4^r(7,10,14,19,26,35,47)$. Thus powers of $4$ do not introduce a new dynamical mechanism; they reveal a scaling symmetry already present in the base system.

The billion-scale base computation gives strong finite evidence for the two-cycle picture, but the global all-integers convergence statement remains a conjecture. The main rigorous contribution of this paper is the exact transfer theorem: once the behavior of the base system is established on any corresponding set of starts, the behavior on its $4^r$-scaled copy follows with no additional dynamical uncertainty.

\begin{thebibliography}{9}
\bibitem{Lagarias1985}
J. C. Lagarias, ``The 3x+1 problem and its generalizations,'' \emph{The American Mathematical Monthly}, 92(1), 3--23, 1985. DOI: 10.2307/2322189.

\bibitem{BanazadehBase}
F. Banazadeh, ``A Ternary (4k + 1(2))/3 Collatz-Type Map with Two Attracting Cycles: Theoretical Analysis and Exhaustive Computational Verification up to $10^9$,'' 2026. Zenodo DOI: 10.5281/ZENODO.22195651.

\bibitem{BanazadehRepo}
F. Banazadeh, ``Sequence Computations,'' GitHub repository, 2026.\\
\url{https://github.com/FarhadBanazadeh/sequence-computations}

\end{thebibliography}

\end{document}
